\subsection{محدودیت‌ها و تحلیل شکست}
تحلیل عملکرد الگوریتم فشرده‌ساز \lr{FRanDI} نقاط حساس و محدودیت‌های زیر را برجسته می‌سازد:
\begin{itemize}
    \item \textbf{افت دقت قبل از بازیابی لورا:} اعمال فاکتورگیری خطی در غیاب کالیبراسیون و آموزش جبرانی کافی لایه‌های \lr{LoRA}، افت دقتی در حدود ۲۰ الی ۳۰ درصد روی انکودرها به جا می‌گذارد که نیاز به دوره‌های فاین‌تیونینگ کوتاه‌مدت پس از فشرده‌سازی را الزامی می‌کند.
    \item \textbf{ناسازگاری ساختاری با انکودر \lr{RoBERTa}:} در مدل \lr{RoBERTa} به دلیل لایه‌بندی متفاوت و حساسیت شدید استخراج بازنمایی‌ها، ساختار فشرده‌ساز در تست غربالگری اولیه به وضعیت \lr{SKIP\_ROBERTA\_ARCHITECTURAL} منتج شده است.
    \item \textbf{هزینه محاسباتی فاکتورگیری چولسکی و حل مثلثی:} در لایه‌های با ابعاد بزرگ ($d_{\text{in}} > 4096$)، تشکیل ماتریس کوواریانس و اعمال $\mathcal{O}(d_{\text{in}}^3)$ تجزیه چولسکی می‌تواند زمان آماده‌سازی اولیه فشرده‌سازی را در مقایسه با روش‌های تنکی ایستا به طور چشمگیری افزایش دهد.
\end{itemize}
